\providecommand{\itescastudentname}{Martin Jonathan de la Cruz}
\providecommand{\itescastudentshort}{M. J. de la Cruz}
\providecommand{\itescastudentid}{ITESCA-ISC}
\providecommand{\itescafaculty}{Instituto Tecnologico Superior de Cajeme}
\providecommand{\itescafacultyname}{Ingenieria en Sistemas Computacionales}
\providecommand{\itescacoursename}{Taller de Sistemas Operativos}
\providecommand{\itescacoursecode}{ISC-TALLER-DE-SISTEMAS-OPERATIVOS}
\providecommand{\itescadepartment}{Taller de Sistemas Operativos}
\providecommand{\itescateachingfigure}{Docente de la materia por definir}
\providecommand{\itescaprogramlevel}{Licenciatura}
\providecommand{\itescaprogramperiod}{Periodo academico por definir}
\providecommand{\itescamodality}{Modalidad por definir}
\providecommand{\itescaprofessionalprofile}{La formacion en Ingenieria en Sistemas Computacionales exige documentar problemas, procesos y decisiones tecnicas con claridad, trazabilidad y sentido de mejora.}
\providecommand{\itescatransfercontext}{Taller de Sistemas Operativos, contexto profesional y aplicacion disciplinar}
\providecommand{\itescalocation}{Cajeme, Sonora}
\providecommand{\itescalocationname}{Cajeme, Sonora}
\providecommand{\itescabibliography}{bibliografia-itesca-isc,ITESCA/ingenieria-en-sistemas-computacionales/taller-de-sistemas-operativos-isc/taller-de-sistemas-operativos}
\providecommand{\itescageneralbib}{bibliografia-itesca-isc.bib}
\providecommand{\itescacoursebib}{taller-de-sistemas-operativos.bib}

\providecommand{\itescamethodologycore}{Problema, evidencia, analisis, postura y transferencia}
\providecommand{\itescadidactictechniques}{Mapa conceptual, diagrama, estudio de caso, cuadro comparativo y proyecto integrador}

\def\itescaactivitytitle {Taller de Sistemas Operativos}
\def\itescaactivitysubtitle {Presentacion base de materia}
\def\itescaactivityweek {Plantilla base}
\def\itescaactivityname {Punto de entrada canonico de materia}
\def\itescasessionname {Taller de Sistemas Operativos}
\def\itescablockname {tronco comun ISC}
\def\itescamainproduct {Presentacion academica base de la materia}
\def\itescadeliverydate {\today}
\def\itescareporttemplatefile {reporte-taller-de-sistemas-operativos-isc.tex}

% ==============================================================================
% PLANTILLA DE PRESENTACION - ITESCA
% Inspirada en la estructura editorial de las presentaciones UnADM del repo,
% adaptada a ITESCA con una identidad propia y reusable para carrera o materia.
%
% Uso recomendado:
% 1. Duplica el archivo contenedor por carrera, materia o actividad.
% 2. Actualiza metadatos, titulo visible, actividad, docente, bibliografia y
%    los bloques editables de la seccion "Plantilla Editable".
% 3. Usa la presentacion como traduccion visual del reporte ITESCA:
%    resumen -> objetivo, introduccion -> problema, desarrollo -> ejes,
%    analisis -> lectura critica, postura -> criterio propio, conclusion -> cierre.
% 4. No pegues la consigna completa en las diapositivas. Convierte la actividad
%    en narrativa, criterio, evidencia y consecuencia profesional.
% 5. Si la actividad exige cuadro, esquema, tabla, mapa, flujo o evidencia,
%    construyelo en LaTeX/TikZ cuando sea viable o inserta una imagen nitida.
% 6. Compilar desde la raiz del repositorio con latexmk-build.ps1.
%
% Criterio institucional ITESCA:
% - Priorizar una voz academica clara, tecnica y orientada a resolucion.
% - Relacionar el tema con la carrera, la materia y el contexto profesional.
% - Evitar definiciones sueltas: cada bloque debe mostrar funcion, limite,
%   implicacion o aplicacion real.
% - Traducir teoria a evidencia, procedimiento, decision o caso de uso.
% - Cerrar con postura, criterio o consecuencia transferible; no con resumen plano.
% - Convertir cada diapositiva en una unidad editorial: pregunta, evidencia,
%   lectura critica y consecuencia.
% - Mantener identidad ITESCA no solo en logotipos, tambien en el enfoque:
%   precision tecnica, pertinencia institucional y transferencia profesional.
%
% Estructura sugerida para cada exposicion:
% [ ] Tecnica didactica declarada.
% [ ] Portada institucional con actividad y datos actualizados.
% [ ] Tesis o pregunta rectora.
% [ ] Objetivo y problema central.
% [ ] Desarrollo en 2-4 ejes claros.
% [ ] Evidencia o producto visual, si la actividad lo requiere.
% [ ] Vinculacion con carrera, proceso o contexto profesional.
% [ ] Postura personal y cierre.
% [ ] Fuentes y bibliografia consistentes con el reporte.
%
% Checklist antes de exponer:
% [ ] Titulo, actividad, asignatura, matricula y docente actualizados.
% [ ] La narrativa responde a la consigna sin copiar instrucciones literales.
% [ ] Cada diapositiva cumple una funcion concreta.
% [ ] Hay criterio propio y no solo recopilacion de conceptos.
% [ ] La presentacion conecta con la carrera o el entorno profesional.
% [ ] El PDF compila sin errores y los recursos institucionales cargan bien.
% ============================================================================== 

% ------------------------------------------------------------------------------
% METADATOS EDITABLES
% ------------------------------------------------------------------------------
\providecommand{\itescastudentname}{Martin Jonathan de la Cruz}
\providecommand{\itescastudentshort}{M. J. de la Cruz}
\providecommand{\itescastudentid}{ITESCA}
\providecommand{\itescauniversityname}{Instituto Tecnologico Superior de Cajeme}
\providecommand{\itescafacultyname}{Programa academico ITESCA}
\providecommand{\itescacoursename}{Nombre de la asignatura}
\providecommand{\itescacoursecode}{ITESCA}
\providecommand{\itescaactivitytitle}{Titulo de la presentacion}
\providecommand{\itescaactivitysubtitle}{Actividad X}
\providecommand{\itescaactivityweek}{Semana X}
\providecommand{\itescateachingfigure}{Nombre por definir}
\providecommand{\itescastudentemail}{correo.institucional@itesca.edu.mx}
\providecommand{\itescaprogramlevel}{Nivel academico por definir}
\providecommand{\itescaprogramperiod}{Periodo academico por definir}
\providecommand{\itescamodality}{Modalidad por definir}
\providecommand{\itescaactivityname}{Actividad por definir}
\providecommand{\itescasessionname}{Sesion por definir}
\providecommand{\itescablockname}{Bloque por definir}
\providecommand{\itescamainproduct}{Producto academico por definir}
\providecommand{\itescaprofessionalprofile}{Perfil profesional por definir}
\providecommand{\itescainstitutionalaxis}{Formacion academica, identidad institucional y transferencia profesional}
\providecommand{\itescaevidencefocus}{Evidencia academica, producto visible y criterio propio}
\providecommand{\itescaqualitycriterion}{Claridad, pertinencia, evidencia, analisis y cierre transferible}
\providecommand{\itescatransfercontext}{Carrera, asignatura, entorno profesional o trayectoria formativa}
\providecommand{\itescadeliverydate}{\today}
\providecommand{\itescalocationname}{Cajeme, Sonora}
\providecommand{\itescageneralbib}{bibliografia-itesca.bib}
\providecommand{\itescacoursebib}{bibliografia de la materia}
\providecommand{\itescareporttemplatefile}{reporte.tex}
\providecommand{\itescaasset}{ITESCA/assets-itesca/itesca-monograma.png}
\providecommand{\itescaofficialasset}{ITESCA/assets-itesca/web/logo-itesca-contorno.png}
\providecommand{\itescasecondaryasset}{ITESCA/assets-itesca/itesca-monograma.png}
\providecommand{\itescacampusurl}{https://www.itesca.edu.mx/}
\providecommand{\itescaportalurl}{https://www.itesca.edu.mx/portalacademico/}
\providecommand{\itescavirtualurl}{https://www.itesca.edu.mx/modulos/oferta/ead/index.php}
\providecommand{\itescabibliotecaurl}{https://www.itesca.edu.mx/modulos/estudiantes/centroinfo.php}

\documentclass[
  spanish,
  aspectratio=169,
  xcolor={dvipsnames,table}
]{beamer}

\geometry{paperwidth=19.2cm,paperheight=10.8cm}

\usepackage[utf8]{inputenc}
\usepackage[T1]{fontenc}
\usepackage[spanish,es-tabla]{babel}
\usepackage[scaled=.96]{helvet}
\renewcommand{\familydefault}{\sfdefault}
\usepackage{booktabs}
\usepackage{graphicx}
\usepackage{ragged2e}
\usepackage{tikz}
\usepackage{hyperref}
\usetikzlibrary{calc,positioning,fit,shadows.blur}

\definecolor{itescaRedDark}{HTML}{8F1117}
\definecolor{itescaRed}{HTML}{D70F18}
\definecolor{itescaCharcoal}{HTML}{141414}
\definecolor{itescaSilver}{HTML}{D9D9D9}
\definecolor{itescaPaper}{HTML}{F7F7F7}
\definecolor{itescaInk}{HTML}{181818}
\definecolor{itescaMuted}{HTML}{6A6A6A}
\colorlet{itescaGreenDark}{itescaRedDark}
\colorlet{itescaGreen}{itescaCharcoal}
\colorlet{itescaGold}{itescaSilver}

\mode<presentation>{
  \usetheme{default}
  \usefonttheme{professionalfonts}
  \setbeamertemplate{navigation symbols}{}
  \setbeamertemplate{blocks}[rounded][shadow=false]
  \setbeamercovered{invisible}
}
\setbeamersize{text margin left=0.55cm,text margin right=0.55cm}

\hypersetup{
  colorlinks=true,
  urlcolor=itescaRedDark,
  linkcolor=itescaRedDark
}

\setbeamercolor{background canvas}{bg=white}
\setbeamercolor{normal text}{fg=itescaInk,bg=white}
\setbeamercolor{structure}{fg=itescaRedDark}
\setbeamercolor{frametitle}{fg=white,bg=itescaCharcoal}
\setbeamercolor{title}{fg=white,bg=itescaCharcoal}
\setbeamercolor{subtitle}{fg=white}
\setbeamercolor{block title}{fg=white,bg=itescaRedDark}
\setbeamercolor{block body}{fg=itescaInk,bg=itescaPaper}
\setbeamercolor{alerted text}{fg=itescaRed}
\setbeamerfont{title}{size=\LARGE,series=\bfseries}
\setbeamerfont{frametitle}{size=\large,series=\bfseries}
\setbeamerfont{block title}{series=\bfseries}
\setbeamertemplate{itemize item}{\textcolor{itescaRedDark}{\large$\blacktriangleright$}}
\setbeamertemplate{itemize subitem}{\textcolor{itescaSilver}{\scriptsize$\blacksquare$}}
\setbeamertemplate{enumerate item}{\textcolor{itescaRedDark}{\insertenumlabel.}}

\setbeamertemplate{frametitle}{
  \nointerlineskip
  \begin{beamercolorbox}[wd=\textwidth,ht=1.02cm,dp=0.22cm,leftskip=0.35cm,rightskip=0.25cm]{frametitle}
    \begin{minipage}[c]{0.72\textwidth}
      \insertframetitle
    \end{minipage}
    \hfill
    \raisebox{-0.07cm}{\includegraphics[height=0.44cm]{\itescaofficialasset}}
  \end{beamercolorbox}
  \vspace{-0.04cm}
  {\color{itescaRed}\rule{\textwidth}{1.3pt}}
}

\setbeamertemplate{footline}{
  \leavevmode%
  \hbox{%
    \begin{beamercolorbox}[wd=.34\paperwidth,ht=2.7ex,dp=1ex,leftskip=1em]{author in head/foot}%
      \color{white}\itescastudentshort
    \end{beamercolorbox}%
    \begin{beamercolorbox}[wd=.40\paperwidth,ht=2.7ex,dp=1ex,center]{title in head/foot}%
      \color{white}\itescacoursename
    \end{beamercolorbox}%
    \begin{beamercolorbox}[wd=.26\paperwidth,ht=2.7ex,dp=1ex,rightskip=1em plus 1fil]{date in head/foot}%
      \color{white}\itescacoursecode\hfill\insertframenumber/\inserttotalframenumber
    \end{beamercolorbox}%
  }%
  \vskip0pt%
}
\setbeamercolor{author in head/foot}{bg=itescaCharcoal}
\setbeamercolor{title in head/foot}{bg=itescaRedDark}
\setbeamercolor{date in head/foot}{bg=itescaCharcoal}

\newcommand{\accentbar}{%
  \begin{tikzpicture}[remember picture,overlay]
    \path[use as bounding box] (0,0) rectangle (0,0);
    \fill[itescaRedDark] (current page.south west) rectangle ([xshift=0.56\paperwidth,yshift=0.08cm]current page.south west);
    \fill[itescaSilver] ([xshift=0.56\paperwidth]current page.south west) rectangle ([xshift=\paperwidth,yshift=0.08cm]current page.south west);
  \end{tikzpicture}
}

\newcommand{\itescasectionframe}[1]{%
  \begin{frame}[plain]
    \begin{tikzpicture}[remember picture,overlay]
      \path[use as bounding box] (0,0) rectangle (0,0);
      \fill[itescaCharcoal] (current page.south west) rectangle (current page.north east);
      \node[opacity=0.08,anchor=east] at ([xshift=-0.35cm]current page.east) {\includegraphics[width=0.42\paperwidth]{\itescaofficialasset}};
      \fill[itescaRedDark] ([xshift=1.05cm,yshift=2.25cm]current page.south west) rectangle ([xshift=1.25cm,yshift=6.40cm]current page.south west);
      \fill[itescaSilver] ([xshift=1.48cm,yshift=2.25cm]current page.south west) rectangle ([xshift=1.68cm,yshift=5.25cm]current page.south west);
      \node[
        anchor=west,
        align=left,
        text=white,
        text width=0.54\paperwidth,
        font=\fontsize{26}{31}\selectfont\bfseries
      ] at ([xshift=2.08cm,yshift=0.55cm]current page.west) {#1};
      \node[
        anchor=west,
        align=left,
        text=white!78,
        font=\large
      ] at ([xshift=2.08cm,yshift=-0.55cm]current page.west) {\itescafacultyname};
    \end{tikzpicture}
  \end{frame}
}

\newcommand{\itescasection}[1]{%
  \section{#1}
  \itescasectionframe{#1}
}

\title[\itescacoursecode]{\itescaactivitytitle}
\subtitle{\itescaactivitysubtitle}
\author[\itescastudentshort]{\texorpdfstring{\itescastudentname\\\footnotesize Matricula: \itescastudentid\\\href{mailto:\itescastudentemail}{\itescastudentemail}}{\itescastudentname, Matricula: \itescastudentid}}
\institute[ITESCA]{\itescauniversityname\\\itescafacultyname\\\itescacoursecode\ \textperiodcentered\ \itescacoursename}
\date[\itescaactivityweek]{\itescalocationname\\\itescaactivityweek\\\itescaprogramperiod\\\itescadeliverydate}

\begin{document}

\begin{frame}[plain]
  \begin{tikzpicture}[remember picture,overlay]
    \path[use as bounding box] (0,0) rectangle (0,0);
    \fill[itescaRedDark,opacity=0.97] (current page.south west) rectangle ([xshift=0.58\paperwidth]current page.north west);
    \fill[itescaCharcoal,opacity=0.95] ([xshift=0.58\paperwidth]current page.south west) rectangle (current page.north east);
    \node[anchor=center,opacity=0.10] at (current page.center) {\includegraphics[width=0.70\paperwidth]{\itescaasset}};
    \node[anchor=north west] at ([xshift=0.58cm,yshift=-0.46cm]current page.north west) {\includegraphics[width=4.10cm]{\itescaofficialasset}};
    \fill[itescaRed] ([xshift=0.60\paperwidth,yshift=1.18cm]current page.south west) rectangle ([xshift=0.93\paperwidth,yshift=1.30cm]current page.south west);
    \fill[white,opacity=0.15] ([xshift=0.60\paperwidth,yshift=0.92cm]current page.south west) rectangle ([xshift=0.84\paperwidth,yshift=1.02cm]current page.south west);
  \end{tikzpicture}

  \vspace*{1.80cm}
  \begin{minipage}{0.55\paperwidth}
    {\color{white}\fontsize{25}{30}\selectfont\bfseries\raggedright \itescaactivitytitle\par}
    \vspace{0.22cm}
    {\color{white!86}\large \itescaactivitysubtitle}\\[0.32cm]
    {\color{itescaSilver}\rule{0.82\linewidth}{1.3pt}}\\[0.42cm]
    {\color{white}\itescastudentname}\\
    {\color{white!82}\footnotesize Matricula: \itescastudentid}\\
    {\color{white!76}\footnotesize \itescastudentemail}
  \end{minipage}
  \hfill
  \begin{minipage}{0.34\paperwidth}
    \raggedleft
    \includegraphics[width=\linewidth]{\itescaofficialasset}\\[0.34cm]
    \includegraphics[width=0.34\linewidth]{\itescasecondaryasset}\\[0.28cm]
    {\color{white}\small \itescafacultyname}\\
    {\color{white!78}\small \itescacoursename}\\[0.45cm]
    {\color{white!74}\footnotesize \itescacoursecode\ \textperiodcentered\ \itescaactivityweek}\\
    {\color{white!74}\footnotesize \itescaprogramlevel\ \textperiodcentered\ \itescamodality}\\
    {\color{white!74}\footnotesize \itescalocationname}\\
    {\color{white!74}\footnotesize \href{\itescacampusurl}{\itescacampusurl}}
  \end{minipage}
\end{frame}

\begin{frame}{Contenido}
  \tableofcontents
\end{frame}

\itescasection{Identidad Academica}

\begin{frame}{Ficha de la actividad}
  \begin{columns}[T]
    \begin{column}{0.40\textwidth}
      \begin{center}
        \includegraphics[width=0.92\linewidth]{\itescaofficialasset}\\[0.28cm]
        {\small\color{itescaMuted}\itescauniversityname}\\
        {\small\color{itescaMuted}\itescafacultyname}
      \end{center}
    \end{column}
    \begin{column}{0.56\textwidth}
      \begin{block}{Datos generales}
        \textbf{Alumno:} \itescastudentname\\
        \textbf{Matricula:} \itescastudentid\\
        \textbf{Correo:} \itescastudentemail\\
        \textbf{Asignatura:} \itescacoursename\\
        \textbf{Codigo:} \itescacoursecode\\
        \textbf{Nivel:} \itescaprogramlevel\\
        \textbf{Actividad:} \itescaactivitysubtitle
      \end{block}
      \begin{block}{Figura docente}
        \itescateachingfigure\\
        \itescaprogramperiod\\
        \itescadeliverydate
      \end{block}
    \end{column}
  \end{columns}
  \accentbar
\end{frame}

\begin{frame}{Contexto academico}
  \begin{columns}[T]
    \begin{column}{0.49\textwidth}
      \begin{block}{Enfoque institucional}
        La presentacion funciona como traduccion visual del reporte academico:
        selecciona el problema, organiza ejes y muestra criterio sin convertir
        las diapositivas en un documento de lectura corrida.
      \end{block}
    \end{column}
    \begin{column}{0.47\textwidth}
      \begin{block}{Ruta de contenido}
        \begin{itemize}
          \item Objetivo y alcance.
          \item Contexto, problema o consigna.
          \item Desarrollo por ejes.
          \item Postura personal.
          \item Conclusion y referencias.
        \end{itemize}
      \end{block}
    \end{column}
  \end{columns}
  \accentbar
\end{frame}

\begin{frame}{Criterio institucional ITESCA}
  \begin{columns}[T]
    \begin{column}{0.48\textwidth}
      \begin{block}{Como debe leerse la actividad}
        No basta con definir conceptos. La presentacion debe mostrar que
        problema atiende el tema, que proceso organiza y que criterio permite
        sostener una respuesta academica o tecnica.
      \end{block}
      \begin{block}{Vinculacion con la carrera}
        Cada actividad conviene aterrizarse a la trayectoria formativa,
        asignatura, procedimiento, caso o decision que tenga sentido dentro del
        programa academico de ITESCA.
      \end{block}
    \end{column}
    \begin{column}{0.48\textwidth}
      \begin{alertblock}{Regla editorial}
        Convertir la consigna en criterio. Formula util: problema
        $\rightarrow$ analisis $\rightarrow$ evidencia $\rightarrow$ consecuencia.
      \end{alertblock}
      \begin{block}{Voz esperada}
        Clara, sobria, argumentada y orientada a aplicacion. Evitar diapositivas
        saturadas, citas sin funcion o definiciones pegadas sin lectura propia.
      \end{block}
    \end{column}
  \end{columns}
  \accentbar
\end{frame}

\begin{frame}{Control editorial de profundidad}
  \begin{columns}[T]
    \begin{column}{0.48\textwidth}
      \begin{block}{Lectura de consigna}
        La actividad debe convertirse en una pregunta de trabajo: que se debe
        comprender, comparar, resolver, evidenciar o justificar.
      \end{block}
      \begin{block}{Evidencia}
        \itescaevidencefocus. La evidencia necesita rotulo, contexto y lectura
        propia para no quedar como adorno.
      \end{block}
    \end{column}
    \begin{column}{0.48\textwidth}
      \begin{block}{Criterio de calidad}
        \itescaqualitycriterion. Cada diapositiva debe aportar una decision:
        delimitar, explicar, demostrar, valorar o cerrar.
      \end{block}
      \begin{alertblock}{Nivel minimo}
        Si hay producto o conclusion, la exposicion debe llegar por lo menos a
        analisis: diferencia, implicacion y postura.
      \end{alertblock}
    \end{column}
  \end{columns}
  \accentbar
\end{frame}

\begin{frame}{Ficha editorial de la exposicion}
  \scriptsize
  \begin{tabular}{p{0.24\textwidth}p{0.66\textwidth}}
    \toprule
    \textbf{Elemento} & \textbf{Criterio} \\
    \midrule
    Programa & \itescafacultyname\\
    Nivel / modalidad & \itescaprogramlevel\ / \itescamodality\\
    Actividad & \itescaactivityname\\
    Sesion / bloque & \itescasessionname\ / \itescablockname\\
    Producto esperado & \itescamainproduct\\
    Eje institucional & \itescainstitutionalaxis\\
    Transferencia & \itescatransfercontext\\
    Perfil profesional & \itescaprofessionalprofile\\
    \bottomrule
  \end{tabular}
  \accentbar
\end{frame}

\itescasection{Enfoque de Contenido}

\begin{frame}{Correspondencia con la plantilla de informe}
  \begin{block}{Logica recomendada}
    La presentacion no debe repetir el reporte completo. Debe condensar su
    estructura para exponer con claridad: objetivo, problema central,
    desarrollo por ejes, producto o evidencia, postura y cierre.
  \end{block}
  \begin{columns}[T]
    \begin{column}{0.48\textwidth}
      \begin{block}{Secciones del reporte}
        \begin{itemize}
          \item Resumen / abstract.
          \item Introduccion.
          \item Desarrollo.
          \item Postura personal.
          \item Conclusion.
          \item Bibliografia.
        \end{itemize}
      \end{block}
    \end{column}
    \begin{column}{0.48\textwidth}
      \begin{block}{Traduccion a diapositivas}
        \begin{itemize}
          \item Objetivo y tesis.
          \item Contexto y problema.
          \item Dos o tres ejes de analisis.
          \item Producto visual o caso.
          \item Lectura critica.
          \item Cierre y fuentes.
        \end{itemize}
      \end{block}
    \end{column}
  \end{columns}
  \accentbar
\end{frame}

\begin{frame}{Narrativa sugerida}
  \begin{enumerate}
    \item Abrir con una afirmacion fuerte y no con una definicion blanda.
    \item Traducir la consigna a una pregunta o problema central.
    \item Desarrollar solo los ejes que sostienen la postura.
    \item Usar producto visual cuando la actividad lo pida.
    \item Cerrar con una implicacion academica, institucional o profesional.
  \end{enumerate}
  \begin{alertblock}{Regla de contenido}
    Cada diapositiva debe responder una pregunta concreta. Si una diapositiva no
    agrega criterio, debe comprimirse o eliminarse.
  \end{alertblock}
  \accentbar
\end{frame}

\begin{frame}{Matriz narrativa de calidad}
  \scriptsize
  \begin{tabular}{p{0.17\textwidth}p{0.33\textwidth}p{0.40\textwidth}}
    \toprule
    \textbf{Momento} & \textbf{Pregunta guia} & \textbf{Salida esperada} \\
    \midrule
    Apertura & Que problema atiende la actividad? & Tesis breve, no definicion generica. \\
    Desarrollo & Que conceptos o procesos explican el problema? & Ejes con funcion, limite e implicacion. \\
    Evidencia & Que producto demuestra el trabajo? & Cuadro, esquema, caso, tabla o procedimiento. \\
    Analisis & Que significa la evidencia? & Lectura propia y consecuencia academica. \\
    Transferencia & Donde se aplica en la carrera? & Decision, habito, herramienta o mejora profesional. \\
    Cierre & Que postura queda sostenida? & Conclusion con criterio y aprendizaje. \\
    \bottomrule
  \end{tabular}
  \accentbar
\end{frame}

\itescasection{Plantilla Editable}

\begin{frame}{Diapositiva tipo: objetivo y problema}
  \begin{block}{Objetivo}
    \textit{Editable:} Analizar, explicar, comparar, interpretar o evaluar el
    tema solicitado por la actividad.
  \end{block}
  \begin{block}{Problema central}
    \textit{Editable:} El problema no se limita a \ldots, sino a \ldots
  \end{block}
  \begin{block}{Tesis de trabajo}
    \textit{Editable:} La posicion que guia esta presentacion es que \ldots
  \end{block}
  \accentbar
\end{frame}

\begin{frame}{Diapositiva tipo: desarrollo por ejes}
  \begin{columns}[T]
    \begin{column}{0.48\textwidth}
      \begin{block}{Eje 1}
        \textit{Concepto} $\rightarrow$ \textit{funcion} $\rightarrow$ \textit{implicacion}.
      \end{block}
      \begin{block}{Eje 2}
        \textit{Problema} $\rightarrow$ \textit{sistema} $\rightarrow$ \textit{consecuencia}.
      \end{block}
    \end{column}
    \begin{column}{0.48\textwidth}
      \begin{block}{Eje 3 o evidencia}
        \textit{Caso, cuadro comparativo, mapa, esquema, tabla o ejemplo aplicado.}
      \end{block}
      \begin{block}{Criterio de seleccion}
        No acumular datos. Mantener solo lo que sostiene la conclusion.
      \end{block}
    \end{column}
  \end{columns}
  \accentbar
\end{frame}

\begin{frame}{Diapositiva tipo: aplicacion institucional y profesional}
  \begin{columns}[T]
    \begin{column}{0.48\textwidth}
      \begin{block}{Vinculacion con la carrera}
        \textit{Editable:} Este tema importa en \ldots porque interviene en
        procesos, decisiones, herramientas o criterios propios de \ldots
      \end{block}
      \begin{block}{Proceso o procedimiento}
        \textit{Editable:} La secuencia relevante es \ldots y su punto critico
        aparece cuando \ldots
      \end{block}
    \end{column}
    \begin{column}{0.48\textwidth}
      \begin{block}{Impacto esperado}
        \textit{Editable:} La aplicacion real del tema se observa en \ldots
      \end{block}
      \begin{alertblock}{Regla}
        No dejar la conclusion en abstracto. Mostrar para que sirve el analisis
        dentro del contexto de ITESCA, la materia o el ejercicio profesional.
      \end{alertblock}
    \end{column}
  \end{columns}
  \accentbar
\end{frame}

\begin{frame}{Diapositiva tipo: criterio de evidencia}
  \begin{columns}[T]
    \begin{column}{0.48\textwidth}
      \begin{block}{Producto visible}
        \textit{Editable:} El producto central de esta actividad es \ldots y
        permite observar \ldots
      \end{block}
      \begin{block}{Criterio de validacion}
        \textit{Editable:} Se considera suficiente cuando muestra \ldots,
        explica \ldots y conecta con \ldots
      \end{block}
    \end{column}
    \begin{column}{0.48\textwidth}
      \begin{block}{Lectura critica}
        \textit{Editable:} La evidencia no solo muestra \ldots; tambien permite
        valorar \ldots
      \end{block}
      \begin{alertblock}{Control}
        Ninguna evidencia debe quedar sin interpretacion: imagen, tabla o
        esquema necesitan una frase de funcion y una consecuencia.
      \end{alertblock}
    \end{column}
  \end{columns}
  \accentbar
\end{frame}

\begin{frame}{Diapositiva tipo: producto visual}
  \begin{block}{Cuando la actividad pida evidencia grafica}
    Insertar aqui un cuadro comparativo, esquema, mapa conceptual, red de ideas,
    linea de tiempo o tabla sintesis. Si el producto se construye fuera de LaTeX,
    usar captura nitida y agregar una frase que explique su funcion.
  \end{block}
  \begin{center}
    \fbox{\parbox[c][4.2cm][c]{0.84\linewidth}{\centering Espacio reservado para el producto visual o evidencia central}}
  \end{center}
  \accentbar
\end{frame}

\begin{frame}{Diapositiva tipo: cierre}
  \begin{block}{Postura personal}
    \textit{Editable:} Desde mi perspectiva, el punto mas importante no es
    \ldots, sino \ldots, porque \ldots
  \end{block}
  \begin{block}{Conclusion}
    \textit{Editable:} En sintesis, el valor del tema esta en \ldots y su
    aplicacion real se observa en \ldots
  \end{block}
  \begin{alertblock}{Cierre util}
    No terminar con resumen plano. Terminar con consecuencia, criterio o
    aprendizaje transferible.
  \end{alertblock}
  \accentbar
\end{frame}

\itescasection{Fuentes y Revision}

\begin{frame}{Fuentes y bibliografia}
  \begin{columns}[T]
    \begin{column}{0.48\textwidth}
      \begin{block}{Bibliografia general}
        \texttt{\itescageneralbib}
      \end{block}
      \begin{block}{Plantilla de informe base}
        \texttt{\itescareporttemplatefile}
      \end{block}
      \begin{block}{Portal Academico}
        \href{\itescaportalurl}{\itescaportalurl}
      \end{block}
    \end{column}
    \begin{column}{0.48\textwidth}
      \begin{block}{Fuentes de la materia}
        \texttt{\itescacoursebib}
      \end{block}
      \begin{block}{Sitio institucional}
        \href{\itescacampusurl}{\itescacampusurl}
      \end{block}
      \begin{block}{ITESCA Virtual / Biblioteca}
        \href{\itescavirtualurl}{ITESCA Virtual}\\
        \href{\itescabibliotecaurl}{Biblioteca institucional}
      \end{block}
    \end{column}
  \end{columns}
  \begin{alertblock}{Regla}
    La presentacion no sustituye la bibliografia del reporte. Solo sintetiza las
    fuentes que sostienen la exposicion.
  \end{alertblock}
  \accentbar
\end{frame}

\begin{frame}{Checklist de revision}
  \begin{itemize}
    \item La portada y los datos de actividad estan actualizados.
    \item La narrativa responde a la planeacion sin copiarla literalmente.
    \item Cada diapositiva tiene una funcion clara.
    \item El contenido mantiene enfoque personal, criterio y postura.
    \item El producto visual aparece cuando la actividad lo exige.
    \item La evidencia esta interpretada y no solo insertada.
    \item La exposicion conecta el tema con la carrera o el entorno profesional.
    \item La identidad ITESCA aparece en datos, enfoque, transferencia y cierre.
    \item La conclusion traduce el tema a una consecuencia practica.
    \item Las fuentes finales coinciden con el reporte y la bibliografia usada.
    \item Si una diapositiva repite el reporte palabra por palabra, se recorta.
  \end{itemize}
  \accentbar
\end{frame}

\end{document}






